\chapter{Conclusion: 다음 학습 축의 조건과 연구 순서}
\label{STC-S17}

본 연구의 답은 단순한 찬반이 아니다. Sleep-time compute은 release-time pre/post-training과
latency-critical test-time learning 사이에, accumulated wake experience를 reusable later-wake state로
바꾸는 deferred lifecycle budget을 추가한다. 이 기능적 축은 이미 external background synthesis에서
product evidence를 갖는다. 그러나 broad per-user foundation-weight sleep이 주류가 되거나 보편 scaling
law를 가진다는 증거는 없다.

\section{처음의 질문에 대한 직접 답}

\textbf{무슨 문제를 푸는가?} Static deployment와 persistent agent experience 사이의 간극, 반복되는 long
context 비용, cross-session strategy learning, freshness/contradiction, finite-memory compaction, wake latency와
learning throughput의 충돌을 푼다. 하나의 method가 이 문제를 모두 풀지는 않는다.

\textbf{Sleep이 최선인가?} One-off document, exact volatile fact, unpredictable query에는 long context,
RAG와 temporal external memory가 더 강하다. Stable high-reuse procedure, multi-episode pattern,
background conflict resolution에서는 deferred transformation이 유망하다. Continual learning에서는 sleep이
replay, regularization, isolation, validation을 실행하는 phase를 제공할 뿐 capacity theorem을 없애지 않는다.

\textbf{Parametric인가 external인가?} Binary choice가 아니다. External source-of-record에 canonical
evidence와 provenance를 보존하고, calibrated reuse와 stability가 충분한 subset만 reversible parametric
artifact로 compile하는 hybrid가 현재 가장 방어적이다. Changing fact, deletion-sensitive episode,
low-reuse item은 external에 남는다.

\textbf{업계와 학계는 어디로 가는가?} 업계는 external synthesis, provenance, control과 background job으로,
학계는 learned memory writer/router, neural fast state, continual-learning mechanism, compaction theory와
benchmarks로 수렴한다. Lifecycle이 mainstream이 되어도 이름은 dreaming, reflection, synthesis 또는
memory로 남을 수 있다.

\textbf{Scaling law는 무엇인가?} 지금은 법칙이 없다. 먼저 complete state/work identity, reuse와 cadence의
restricted break-even, capacity knee와 state-migration crossover를 측정한다. 그다음 evidence quality,
valid reuse, cycle count, recursive depth, medium, hardware를 교차한 held-out predictive surface가 반복되면
비로소 제한된 scaling law를 말할 수 있다.

\textbf{Memory-device 회사의 기회는 무엇인가?} 특정 parametric sleep adoption에 베팅하기보다 versioned
derived state, provenance, temporal index, snapshot/rollback, heterogeneous state cache, state-local processing,
lifecycle telemetry의 실제 trace를 공동 측정해야 한다. 상세 제품 hypothesis는 patent-first review 뒤에
공개한다.

\section{권장 연구 순서}

\begin{enumerate}
  \item \textbf{External lifecycle baseline}: versioned raw/summary/vector/graph와 exact delete/rollback을 먼저
  구현해 source-of-truth와 metric을 고정한다.
  \item \textbf{Matched benchmark}: 동일 event/query/resource에서 raw, compressed, latent, LoRA, static
  mixture, past-only router를 128+ cycles 비교한다.
  \item \textbf{Reuse/capacity sweep}: $R$, volatility, active capacity, recursive depth, cadence를 독립적으로
  변화시켜 crossover와 knee를 찾는다.
  \item \textbf{Selective promotion}: stable procedural subset에만 versioned adapter/latent compile을 허용하고
  rebase, demotion, delete fault를 주입한다.
  \item \textbf{Trace-driven infra DSE}: shared, time-partitioned, separate, shard-local placement에서 base
  residency와 모든 state-migration bytes를 계측한다.
  \item \textbf{Device co-design}: public trace로 capacity, endurance, tail latency, COW, secure reclamation의
  sensitivity를 검증한 뒤 product research로 승격한다.
\end{enumerate}

\section{마지막 framing}

앞으로의 학습을 세 phase로 보는 framing은 유용하다.

\begin{align*}
\underbrace{\text{pre/post-training}}_{\text{release 전 shared capability}}
&\rightarrow
\underbrace{\text{test-time learning}}_{\text{wake adaptation}}\\
&\rightarrow
\underbrace{\text{sleep-time learning}}_{\text{deferred consolidation}}
\rightarrow
\underbrace{\text{versioned long-term memory}}_{\text{later-wake reuse}}.
\end{align*}

다만 화살표는 일방향이 아니다. Delete와 correction은 long-term memory에서 derived state로 역전파되고,
base release는 adapter와 index를 재검증하며, wake failure는 다음 sleep candidate를 만든다. 미래의 scaling
object는 checkpoint 하나가 아니라 여러 clock과 memory medium, provenance와 transaction을 가진 system일
가능성이 높다.

\begin{keypoint}[최종 verdict]
Sleep-time compute은 promising한 새 lifecycle 축이지만, 현재의 승자는 broad weight update가 아니라
external background consolidation이다. 가장 유망한 다음 단계는 external source-of-record 위의 selective,
reversible promotion이며, 그 가치는 sleep FLOPs가 아니라 valid later-wake reuse per complete lifecycle
resource로 검증해야 한다.
\end{keypoint}
